\chapter{Introduction}
\label{ch:intro}

\section{The question behind the project}

In the television series \emph{Person of Interest}~\citep{poi2011}, a system
called the Machine watches the cameras of a city. It recognises people,
follows them from one camera to the next and predicts which of them will be
involved in violence. As fiction it combines several technologies that are
real today and several that are not. Its lowest layer is the one this
project is about: a perception system that looks at camera feeds and says,
reliably and in real time, what objects are present and where they are.

The supervised project topic, \emph{YOLO-Based Cross Feature Fusion Mamba
Net for Object Detection}, names three ingredients of such a layer. YOLO is
the family of single-stage detectors that made real-time detection practical
on ordinary hardware~\citep{yolo26_2026,wang2024yolov10,tian2025yolov12}.
Cross-feature fusion is the combination of information from different
sources, such as a visible-light camera and a thermal camera, or features
taken at different image scales. Mamba is a selective state-space model,
a sequence model whose cost grows linearly with sequence length and which
summarises everything it has seen in a fixed-size hidden
state~\citep{gu2023mamba}.

This dossier sets out how these ingredients can be combined into work that is
new, measurable and publishable, and how the result can be built in phases
that end with a system running on live cameras.

\section{From a fictional Machine to a defensible research problem}

A research contribution needs a precise problem, a measurable claim and an
honest account of prior work. The Machine is none of these things. Three
choices narrow it into something that can be studied.

The first choice is the scene. Objects that occupy few pixels are the hard
case for every detector. A person 200 metres from a camera or a car seen from
a drone may cover fewer than $16\times16$ pixels. Small-object detection
remains the weakest part of the COCO benchmark~\citep{lin2014coco}. It has
dedicated benchmarks, including VisDrone~\citep{zhu2021visdrone},
AI-TOD~\citep{xu2022aitodv2}, SODA~\citep{cheng2023soda} and
RGBT-Tiny~\citep{ying2025rgbttiny}.

The second choice is the sensor. Visible cameras fail at night, in glare
and in fog. Thermal cameras see heat regardless of light but lose colour
and texture. A detector that fuses both is useful exactly when a single
sensor is not~\citep{li2025cfmw,dong2025fusionmamba}.

The third choice is time. A camera does not produce isolated photographs.
It produces a stream in which a faint object seen in one frame is usually
still there in the next. Most published detectors, including all four
source papers studied here, process each image independently. State-space
models were designed for sequences, so carrying their memory across frames
is a natural use of them. As Chapter~\ref{ch:novelty} shows, it is also the
least explored one.

\section{Data in, detections out}
\label{sec:datain}

A common early question is whether the system works only on a dataset that is
handed to it, or whether it can also watch a live camera. It can do both,
and the distinction is worth stating precisely because it organises the
whole plan.

\emph{Training} always uses curated datasets. These are collections of images
or video frames in which every object of interest has been labelled with a
box and a class. The model learns from these labels, and its accuracy is
measured on a held-out portion of the same datasets. This part of the work is
the same whatever the final application.

\emph{Inference}, the use of the trained model, can draw on any source of
images. The same weights can read a folder of test images, an uploaded video
file or a live stream from a webcam, an IP camera or a thermal camera. The
detector itself does not change. What changes is the work around it. Video
must be decoded and its frames kept in order. A live stream must be read
fast enough to keep up, so stale frames must be dropped. A streaming model
must also keep its memory from one frame to the next.
Figure~\ref{fig:glance} shows the four phases that follow from this.

\begin{figure}[p]
\centering
\begin{tikzpicture}[
    col/.style={base, text width=2.5cm, minimum height=1.42cm},
  ]
  % columns (phases) and rows
  \def\xa{0}\def\xb{3.0}\def\xc{6.0}\def\xd{9.0}
  \def\yh{0.15}\def\ys{-1.3}\def\yi{-3.25}\def\ym{-5.35}\def\yp{-7.45}\def\yo{-9.4}
  \def\xt{-3.0}

  % row labels on the far left
  \node[grp, anchor=east] at (\xt+1.05,\ys) {};

  % phase headings
  \node[grp] at (\xa,\yh) {Phase 1\\[-1pt]still images};
  \node[grp] at (\xb,\yh) {Phase 2\\[-1pt]uploaded video};
  \node[grp] at (\xc,\yh) {Phase 3\\[-1pt]live streams};
  \node[grp] at (\xd,\yh) {Phase 4\\[-1pt]multi-camera};
  \node[grp] at (\xt,\yh) {Offline};

  % sources
  \node[col, dat] (s1) at (\xa,\ys) {Dataset folders\\ labelled visible--\\thermal pairs};
  \node[col, dat] (s2) at (\xb,\ys) {Video files\\ single or paired\\ visible--thermal};
  \node[col, dat] (s3) at (\xc,\ys) {Live cameras\\ webcam, RTSP,\\ thermal sensor};
  \node[col, dat] (s4) at (\xd,\ys) {Camera network\\ several synced\\ live streams};

  % ingest
  \node[col, n] (i1) at (\xa,\yi) {Paired loader\\ letterbox and\\ shared geometry};
  \node[col, n] (i2) at (\xb,\yi) {Decoder\\ PyAV or OpenCV,\\ timestamp pairing};
  \node[col, n] (i3) at (\xc,\yi) {Capture threads\\ latest frame wins,\\ queue of one};
  \node[col, n] (i4) at (\xd,\yi) {Camera workers\\ one per stream,\\ shared clock};

  % model row: one network, two modes
  \node[base, fuse, text width=2.5cm, minimum height=1.6cm] (m1) at (\xa,\ym)
       {\textbf{\method}\\ image mode, state\\ reset each frame};
  \node[base, tim, text width=8.5cm, minimum height=1.6cm] (m2) at ({(\xb+\xd)/2},\ym)
       {\textbf{\methodS}\enspace stream mode\\[1pt]
        identical weights, plus the Temporal State Memory\\
        carried from frame to frame at constant cost};

  % post-processing
  \node[col, n] (p1) at (\xa,\yp) {Detections\\ boxes, classes\\ and scores};
  \node[col, n] (p2) at (\xb,\yp) {Tracker\\ BoT-SORT or\\ ByteTrack};
  \node[col, n] (p3) at (\xc,\yp) {Tracker, then\\ overlay renderer\\ and encoder};
  \node[col, n] (p4) at (\xd,\yp) {Anonymous\\ cross-camera\\ association};

  % outputs
  \node[col, fuse] (o1) at (\xa,\yo) {Benchmark scores\\ AP, AP\textsubscript{S},\\ tiny-object AP};
  \node[col, fuse] (o2) at (\xb,\yo) {Annotated video\\ HOTA, IDF1,\\ frame latency};
  \node[col, fuse] (o3) at (\xc,\yo) {Live dashboard\\ under 150\,ms\\ end to end};
  \node[col, fuse] (o4) at (\xd,\yo) {Counts, zones,\\ events and\\ heat maps};

  % arrows down each column
  \foreach \c in {1,2,3,4}{
    \draw[arr] (s\c) -- (i\c);
    \draw[arr] (p\c) -- (o\c);
  }
  \draw[arr] (i1) -- (m1);
  \draw[arr] (m1) -- (p1);
  \foreach \c/\x in {2/\xb,3/\xc,4/\xd}{
    \draw[arr] (i\c.south) -- (\x,{\ym+0.8});
    \draw[arr] (\x,{\ym-0.8}) -- (p\c.north);
  }

  % offline training lane
  \node[base, dat, text width=1.95cm, minimum height=1.6cm] (tr) at (\xt,\ym)
       {Training\\ labelled data,\\ Kaggle T4};
  \node[base, n, text width=1.95cm, minimum height=1.42cm] (ck) at (\xt,\yi)
       {Checkpoints\\ seeds, logs,\\ exports};
  \node[base, n, text width=1.95cm, minimum height=1.42cm] (an) at (\xt,\yo)
       {Error analysis\\ missed and\\ false objects};
  \draw[arr] (tr.east) -- node[lbl, above=0.5pt] {weights} (m1.west);
  \draw[arr] (tr) -- (ck);
  \draw[arr] (o1.west) -- (an.east);
  \draw[darr] (an) -- node[lbl, fill=white, text width=1.6cm, pos=0.5] {new or\\ relabelled\\ data} (tr);
\end{tikzpicture}

\vspace{4mm}
\caption[The system at a glance]{The system at a glance. Training always
uses labelled datasets. The trained weights then serve four kinds of input,
one per phase. Phase 1 evaluates the image-mode network on still images.
Phases 2 to 4 run the same weights in stream mode, where the Temporal State
Memory of \methodS{} carries information between frames. Errors found on
benchmarks and in the field are analysed and turned into new or corrected
training data.}
\label{fig:glance}
\end{figure}


\section{Research questions}

The plan is organised around four questions. Each can be answered by
experiment.

\begin{enumerate}[label=\textsc{rq}\arabic*., leftmargin=3.2em]
  \item Does a state-space memory carried causally across frames improve the
        detection of very small objects in visible--thermal video, compared
        with the same detector applied frame by frame?
  \item Is the selective state-space formulation the right form for that
        memory? Specifically, does it beat a convolutional recurrent cell
        and windowed temporal attention at matched computation?
  \item Can the discretisation step of the state-space model act as a learned
        reliability gate? If so, the detector should degrade gracefully when
        one sensor drops out, drifts out of alignment or delivers frames
        irregularly.
  \item Can the resulting detector run on live streams within a real-time
        budget, on a desktop GPU and on an embedded Jetson module, without
        losing the accuracy measured offline?
\end{enumerate}

\section{What this dossier contributes}

The dossier is a planning document, not a results paper. It contributes the
following.

\begin{itemize}
  \item A critical reading of the four source papers. The reading
        reconstructs their tables from the original PDFs and identifies
        where their gains come from (Chapter~\ref{ch:sources}).
  \item A survey of the state of the art as of October 2026, newest first,
        in six areas. These are real-time detectors, vision state-space
        models, small objects, visible--thermal fusion, video and tracking,
        and foundation models (Chapter~\ref{ch:sota}).
  \item A novelty analysis of 64 closely related papers. It shows which
        directions are saturated and which remain open
        (Chapter~\ref{ch:novelty}).
  \item The design of \method{} and \methodS{}, with equations, complexity
        analysis and falsifiable hypotheses (Chapter~\ref{ch:method}).
  \item A dataset stack, an experimental protocol that meets the standards
        of the target journals, a four-phase roadmap and a notebook workflow
        (Chapters~\ref{ch:data} to~\ref{ch:notebooks}).
  \item An account of the legal and ethical limits that separate this
        system from the fictional Machine (Chapter~\ref{ch:ethics}).
\end{itemize}

\section{How the programme fits together}

Figure~\ref{fig:programme} shows the order of work and how later stages feed
back into earlier ones. It also serves as a reading guide.

\begin{figure}[htbp]
\centering
\begin{tikzpicture}[
    st/.style={base, text width=2.42cm, minimum height=1.55cm},
  ]
  \def\dx{2.92}
  % top row, left to right
  \node[st, n]     (lit)  at (0*\dx,0) {\textbf{Evidence}\\ source papers,\\ state of the art\\ \textit{Ch.~3--4}};
  \node[st, alert] (gap)  at (1*\dx,0) {\textbf{Gap}\\ saturated and\\ open directions\\ \textit{Ch.~5}};
  \node[st, fuse]  (meth) at (2*\dx,0) {\textbf{Method}\\ \method, \methodS,\\ hypotheses\\ \textit{Ch.~6}};
  \node[st, dat]   (data) at (3*\dx,0) {\textbf{Data}\\ core stack,\\ alignment audit\\ \textit{Ch.~7}};
  \node[st, n]     (prot) at (4*\dx,0) {\textbf{Protocol}\\ metrics, baselines,\\ ablations\\ \textit{Ch.~8}};

  % bottom row, right to left
  \node[st, tim]   (phas) at (4*\dx,-2.75) {\textbf{Build}\\ Phases 0--4,\\ notebooks\\ \textit{Ch.~9--10}};
  \node[st, fuse]  (res)  at (3*\dx,-2.75) {\textbf{Results}\\ tables built\\ from run logs\\ \textit{Notebook 15}};
  \node[st, n]     (eth)  at (2*\dx,-2.75) {\textbf{Ethics}\\ privacy review,\\ approvals\\ \textit{Ch.~11}};
  \node[st, alert] (pap)  at (1*\dx,-2.75) {\textbf{Manuscript}\\ journal choice,\\ checklist\\ \textit{Ch.~12}};

  \draw[arr] (lit) -- (gap);
  \draw[arr] (gap) -- (meth);
  \draw[arr] (meth) -- (data);
  \draw[arr] (data) -- (prot);
  \draw[arr] (prot) -- (phas);
  \draw[arr] (phas) -- (res);
  \draw[arr] (res) -- (eth);
  \draw[arr] (eth) -- (pap);

  % feedback loops
  \draw[darr] (res.north west) -- node[lbl, fill=white, pos=0.5] {revise} (meth.south east);
  \draw[darr] (pap.north) -- node[lbl, fill=white, pos=0.5] {re-check} (gap.south);
\end{tikzpicture}

\vspace{2mm}
\caption[The research programme and its feedback loops]{The research
programme. Solid arrows give the order of work and the chapters that
describe each stage. The two dashed arrows are the feedback loops that
matter most. Ablation results revise the method. Prior art is checked again
immediately before submission, because the field publishes new Mamba
detectors every month.}
\label{fig:programme}
\end{figure}


\section{Conventions}

Numbers quoted from other work are taken from the cited paper or its official
repository. The dagger (\unv) marks figures that could only be confirmed from
a secondary source. Dates in tables are the month a work first appeared,
usually as an arXiv preprint. Mean average precision at an IoU threshold of
0.5 is written AP\textsubscript{50}. The COCO average over thresholds 0.5 to
0.95 is written AP or mAP\textsubscript{50:95}. Small-object metrics follow
the definitions in Section~\ref{sec:small}.
